\documentclass[10pt]{article}

\usepackage[margin=1in]{geometry}
\usepackage{amsmath, amssymb}
\usepackage{graphicx}
\usepackage{booktabs}
\usepackage{hyperref}
\usepackage{xcolor}

\title{Selective Predictive State Models: Synthetic Baseline Note}
\author{Malo de Pastor}
\date{\today}

\begin{document}
\maketitle

\begin{abstract}
This note documents a controlled synthetic proof of concept for Selective Predictive State Models. The goal is to test whether a predictive latent model can estimate the reliability of its own cheap future-state prediction and use this estimate to decide when optional refinement compute is worth paying for. The experiment is deliberately synthetic: it validates the mechanism before moving to real visual features and public datasets.
\end{abstract}

\section{Research question}
A latent world model usually learns a transition of the form
\begin{equation}
    z_t, a_t \mapsto \hat z_{t+1},
\end{equation}
where \(z_t\) is the current latent state and \(a_t\) is an action or transition context. This project asks whether such a model can also learn when its own prediction is likely to be unreliable, and whether that reliability estimate can drive compute-aware selective refinement.

\section{Synthetic setup}
Each synthetic sample contains a current latent state \(z_t\), an action \(a_t\), and a future latent state \(z_{t+1}\). Some transitions are marked as predictably difficult before observing the future; these correspond to \emph{expected unreliability}. Some samples also receive a post-hoc future mismatch, representing an \emph{observed surprise}.

The model first produces a cheap prediction \(\hat z_{t+1}^{\mathrm{cheap}}\). A reliability head predicts an unreliability score
\begin{equation}
    r_t = R(z_t, a_t, \hat z_{t+1}^{\mathrm{cheap}}, \phi(a_t)),
\end{equation}
where \(\phi(a_t)\) denotes simple action-derived features such as absolute values, squared values, and norms.

The selector activates optional refinement when
\begin{equation}
    r_t > \tau.
\end{equation}
The final prediction is cheap if the selector does not activate refinement, and refined otherwise.

\section{Diagnostics}
We distinguish four diagnostics:
\begin{center}
\begin{tabular}{llll}
\toprule
Metric & Score & Label & Meaning \\
\midrule
Expected learned & learned unreliability & expected hard transition & anticipates predictable difficulty \\
Expected residual & realized residual & expected hard transition & oracle-style check \\
Observed learned & learned unreliability & observed surprise & tests anticipability of surprise \\
Observed residual & realized residual & observed surprise & detects surprise after observation \\
\bottomrule
\end{tabular}
\end{center}

The realized residual is
\begin{equation}
    e_{t+1} = \lVert \hat z_{t+1} - z_{t+1} \rVert^2.
\end{equation}
It acts as a post-observation surprise score.

\section{Utility}
The compute-aware policy is evaluated with
\begin{equation}
    U = -\operatorname{mean}(e) - \lambda \operatorname{mean}(c),
\end{equation}
where \(e\) is prediction error, \(c\) is compute cost, and \(\lambda\) controls the compute penalty. A successful selective policy should approach all-expensive error while using substantially less compute.

\section{Current status}
The current synthetic baseline validates the mechanism qualitatively. The learned reliability score predicts expected unreliability, residual error detects observed surprise, and an adaptive threshold policy improves the utility--compute trade-off relative to cheap-only and all-expensive policies. The next step is to replace synthetic latents with frozen visual features extracted from real image or video sequences.

\end{document}
